\section{Conclusion}
\label{sec:conclusion}

Verify, then vote improves majority accuracy when agents cite passages that a checker can assess. Verification supplies the main signal: the simpler plurality selector matches the scored rule's accuracy, and the passage controls remove most corrections. Keeping a verified default reduces harmful overturns, but the original scored rule misses its pre-specified harm criterion. The second-pool result supports transfer across agents on shared questions; unseen-question generalization and independent human validation of the failure attribution remain open.

The formal guarantee makes failures inspectable, not necessarily rare. The practical contribution is an explicit selection policy whose passage evidence can be audited and reused for later decisions over existing sealed inputs. A strong judge remains more accurate and cheaper for a first decision; verify-then-vote is useful when this policy control and evidence reuse matter. The artifact and \suppref{S7} make the measured tradeoffs available for re-analysis.
